% Clase 29 — no toda onda electromagnética es luz: el visible es una franja
% estrecha (aproximadamente 400 a 700 nm), con el infrarrojo a longitudes de
% onda mayores (frecuencias menores que el rojo) y el ultravioleta a menores.
% Se marcan también las bandas de la Clase 26. Escala logarítmica en metros.
\begin{tikzpicture}
  \begin{axis}[
    emfig, width=12.5cm, height=4.2cm,
    xmode=log, log basis x=10,
    xmin=1e-9, xmax=1e4, ymin=0, ymax=1,
    ytick=\empty, grid=none,
    xlabel={longitud de onda $\lambda$ (m)},
  ]
    \fill[figamber!45] (axis cs:4e-7,0.15) rectangle (axis cs:7e-7,0.85);
    \fill[figred!18] (axis cs:7e-7,0.15) rectangle (axis cs:1e-3,0.85);
    \fill[figblue!18] (axis cs:1e-8,0.15) rectangle (axis cs:4e-7,0.85);
    \fill[figgray!18] (axis cs:1e-3,0.15) rectangle (axis cs:0.3,0.85);
    \fill[figgray!30] (axis cs:0.3,0.15) rectangle (axis cs:1e4,0.85);
    \node[font=\footnotesize] at (axis cs:6e-8,0.5) {UV};
    \node[font=\footnotesize, rotate=90] at (axis cs:5.3e-7,0.5) {visible};
    \node[font=\footnotesize] at (axis cs:2.6e-5,0.5) {infrarrojo};
    \node[font=\footnotesize] at (axis cs:1.7e-2,0.5) {microondas};
    \node[font=\footnotesize] at (axis cs:60,0.5) {radio (FM, AM)};
  \end{axis}
\end{tikzpicture}
